\subsection*{Dedekind Cuts}

\begin{definitionbox}{Definition}
\begin{definition}[Dedekind Cut]
\label{def:dedekind-cut}
A \emph{Dedekind cut} is a set $\alpha\subseteq\mathbb{Q}$ satisfying:
\[
\alpha\ne\emptyset,\qquad \alpha\ne\mathbb{Q},
\]
\[
p\in\alpha\land q<p \Longrightarrow q\in\alpha,
\]
and
\[
p\in\alpha \Longrightarrow \exists r\in\alpha\,(p<r).
\]
\end{definition}
\end{definitionbox}

\begin{remark*}[Standard quantified statement]
\[
\operatorname{DedekindCut}(\alpha)
\Longleftrightarrow
\alpha\subseteq\mathbb{Q}\land \alpha\ne\emptyset\land \alpha\ne\mathbb{Q}
\land
\forall p,q\in\mathbb{Q}\,((p\in\alpha\land q<p)\Longrightarrow q\in\alpha)
\land
\forall p\in\mathbb{Q}\,(p\in\alpha\Longrightarrow \exists r\in\mathbb{Q}\,(r\in\alpha\land p<r)).
\]
\end{remark*}

\begin{remark*}[Predicate reading]
\[
\operatorname{DedekindCut}(\alpha)
\]
means that $\alpha$ is a Dedekind cut of $\mathbb{Q}$.
\end{remark*}

\begin{remark*}[Interpretation]
A Dedekind cut is a nontrivial downward-closed subset of the rationals with no greatest element.
\end{remark*}

\begin{dependencies}
\begin{itemize}
  \item \hyperref[thm:density-of-q]{Density of $\mathbb{Q}$}
\end{itemize}
\end{dependencies}

\begin{definitionbox}{Definition}
\begin{definition}[The Real Numbers]
\label{def:reals}
The real numbers are the set
\[
\mathbb{R}:=\{\alpha:\alpha\text{ is a Dedekind cut in }\mathbb{Q}\}.
\]
\end{definition}
\end{definitionbox}

\begin{remark*}[Standard quantified statement]
\[
\mathbb{R}=\{\alpha\subseteq\mathbb{Q}:\operatorname{DedekindCut}(\alpha)\}.
\]
\end{remark*}

\begin{remark*}[Interpretation]
The real numbers are constructed as the collection of all Dedekind cuts of rational numbers.
\end{remark*}

\begin{dependencies}
\begin{itemize}
  \item \hyperref[def:dedekind-cut]{Dedekind Cut}
\end{itemize}
\end{dependencies}

\subsection*{Rational Cuts}

\begin{definitionbox}{Definition}
\begin{definition}[Rational Cut]
\label{def:rational-cut}
For $q\in\mathbb{Q}$, define the rational cut
\[
q^{\ast}:=\{r\in\mathbb{Q}:r<q\}.
\]
\end{definition}
\end{definitionbox}

\begin{remark*}[Standard quantified statement]
\[
\forall q\in\mathbb{Q},\qquad
q^{\ast}=\{r\in\mathbb{Q}:r<q\}.
\]
\end{remark*}

\begin{remark*}[Interpretation]
A rational number embeds into the real construction as the cut of all rationals below it.
\end{remark*}

\begin{dependencies}
\begin{itemize}
  \item \hyperref[def:rationals]{Rational Numbers}
  \item \hyperref[def:order-on-q]{Order on $\mathbb{Q}$}
\end{itemize}
\end{dependencies}

\begin{lemmabox}{Lemma}
\begin{lemma}[Rational Cuts Are Cuts]
\label{lem:rational-cut-is-a-cut}
For every $q\in\mathbb{Q}$, the set $q^{\ast}$ is a Dedekind cut.
\hyperref[prf:rational-cut-is-a-cut]{\textit{Go to proof.}}
\end{lemma}
\end{lemmabox}

\begin{remark*}[Standard quantified statement]
\[
\forall q\in\mathbb{Q},\qquad \operatorname{DedekindCut}(q^{\ast}).
\]
\end{remark*}

\begin{remark*}[Interpretation]
Every rational cut satisfies the defining cut conditions, so each rational determines a real number.
\end{remark*}

\begin{dependencies}
\begin{itemize}
  \item \hyperref[thm:density-of-q]{Density of $\mathbb{Q}$}
  \item \hyperref[prop:no-least-positive-rational]{No Least Positive Rational}
\end{itemize}
\end{dependencies}

\subsection*{\texorpdfstring{Zero and One in $\mathbb{R}$}{Zero and One in R}}

\begin{definitionbox}{Definition}
\begin{definition}[Zero in $\mathbb{R}$]
\label{def:zero-in-r}
Define
\[
0_{\mathbb{R}}:=(0_{\mathbb{Q}})^{\ast}
=\{r\in\mathbb{Q}:r<0_{\mathbb{Q}}\}.
\]
\end{definition}
\end{definitionbox}

\begin{remark*}[Standard quantified statement]
\[
0_{\mathbb{R}}=(0_{\mathbb{Q}})^{\ast}
=\{r\in\mathbb{Q}:r<0_{\mathbb{Q}}\}.
\]
\end{remark*}

\begin{remark*}[Interpretation]
The real zero is the rational cut determined by rational zero.
\end{remark*}

\begin{dependencies}
\begin{itemize}
  \item \hyperref[def:rational-cut]{Rational Cut}
  \item \hyperref[def:zero-in-q]{Zero in $\mathbb{Q}$}
\end{itemize}
\end{dependencies}

\begin{definitionbox}{Definition}
\begin{definition}[One in $\mathbb{R}$]
\label{def:one-in-r}
Define
\[
1_{\mathbb{R}}:=(1_{\mathbb{Q}})^{\ast}
=\{r\in\mathbb{Q}:r<1_{\mathbb{Q}}\}.
\]
\end{definition}
\end{definitionbox}

\begin{remark*}[Standard quantified statement]
\[
1_{\mathbb{R}}=(1_{\mathbb{Q}})^{\ast}
=\{r\in\mathbb{Q}:r<1_{\mathbb{Q}}\}.
\]
\end{remark*}

\begin{remark*}[Interpretation]
The real one is the rational cut determined by rational one.
\end{remark*}

\begin{dependencies}
\begin{itemize}
  \item \hyperref[def:rational-cut]{Rational Cut}
  \item \hyperref[def:one-in-q]{One in $\mathbb{Q}$}
\end{itemize}
\end{dependencies}

\begin{theorembox}{Theorem}
\begin{theorem}[Zero and One Are Reals]
\label{thm:zero-one-are-reals}
The cuts $0_{\mathbb{R}}$ and $1_{\mathbb{R}}$ are elements of $\mathbb{R}$.
\hyperref[prf:zero-one-are-reals]{\textit{Go to proof.}}
\end{theorem}
\end{theorembox}

\begin{remark*}[Standard quantified statement]
\[
0_{\mathbb{R}}\in\mathbb{R}\land 1_{\mathbb{R}}\in\mathbb{R}.
\]
\end{remark*}

\begin{remark*}[Interpretation]
The designated zero and one cuts are valid real numbers.
\end{remark*}

\begin{dependencies}
\begin{itemize}
  \item
 \hyperref[def:reals]{The Real Numbers}
  \item \hyperref[def:zero-in-r]{Zero in $\mathbb{R}$}
  \item \hyperref[def:one-in-r]{One in $\mathbb{R}$}
  \item \hyperref[lem:rational-cut-is-a-cut]{Rational Cuts Are Cuts}
\end{itemize}
\end{dependencies}

\begin{theorembox}{Theorem}
\begin{theorem}[Zero Is Not One in $\mathbb{R}$]
\label{thm:zero-not-one-in-r}
In $\mathbb{R}$,
\[
0_{\mathbb{R}}\ne 1_{\mathbb{R}}.
\]
\hyperref[prf:zero-not-one-in-r]{\textit{Go to proof.}}
\end{theorem}
\end{theorembox}

\begin{remark*}[Standard quantified statement]
\[
0_{\mathbb{R}}\ne 1_{\mathbb{R}}.
\]
\end{remark*}

\begin{remark*}[Interpretation]
The real constants inherited from rational zero and one remain distinct.
\end{remark*}

\begin{dependencies}
\begin{itemize}
  \item \hyperref[def:zero-in-r]{Zero in $\mathbb{R}$}
  \item \hyperref[def:one-in-r]{One in $\mathbb{R}$}
  \item \hyperref[cor:one-positive-on-q]{One Is Positive in $\mathbb{Q}$}
\end{itemize}
\end{dependencies}
